\subsection*{Questão 05}

\textit{Qual a relação de período e frequência? }

\textbf{Resposta: } 
A relação entre eles é proporcionalidade inversa. Em termos simples: quanto mais rápido algo se repete, menor é o tempo que leva para completar uma única repetição.
\begin{itemize}
    \item Período (T): É o intervalo de tempo necessário para que um ciclo completo ocorra. Sua unidade no Sistema Internacional (SI) é o segundo (s).
    \item Frequência ($f$): É o número de ciclos que ocorrem em uma determinada unidade de tempo (geralmente 1 segundo). Sua unidade no SI é o Hertz (Hz), onde $1\text{ Hz} = 1\text{ ciclo/segundo}$.
\end{itemize}
\textbf{Formúla Matamática:}\par 
$$f = \frac{1}{T} \qquad ou \qquad T = \frac{1}{f}$$